\documentclass[11pt, a4paper]{article}

% --- Minimal Packages ---
\usepackage[utf8]{inputenc}
\usepackage[T1]{fontenc}
\usepackage{geometry}
\geometry{top=3cm, bottom=3cm, left=2.5cm, right=2.5cm}
\usepackage{amsmath, amsthm, amssymb, mathtools}
\usepackage{hyperref}
\usepackage{enumitem}

% --- Macros (Key Objects) ---
\newcommand{\layer}[1]{\mathsf{Layer}_{#1}}
\newcommand{\minLayer}{\mathsf{minLayer}}
\newcommand{\rank}{\rho}
\newcommand{\Tower}{\mathcal{T}}
\newcommand{\Filtration}{\mathbb{F}}
\newcommand{\Process}[1]{(#1_n)_{n \in \mathbb{N}}}
\newcommand{\CondExp}[2]{\mathbb{E}[#1 \mid \mathcal{F}_{#2}]}
\newcommand{\Nat}{\mathbb{N}}
\newcommand{\Real}{\mathbb{R}}
\newcommand{\Prob}{\mathbb{P}}
\newcommand{\Ex}{\mathbb{E}}
\newcommand{\Ind}[1]{\mathbf{1}_{#1}}

% --- Theorem Environments ---
\theoremstyle{plain}
\newtheorem{theorem}{Theorem}[section]
\newtheorem{lemma}[theorem]{Lemma}
\newtheorem{proposition}[theorem]{Proposition}
\newtheorem{corollary}[theorem]{Corollary}

\theoremstyle{definition}
\newtheorem{definition}[theorem]{Definition}
\newtheorem{example}[theorem]{Example}

\theoremstyle{remark}
\newtheorem{remark}[theorem]{Remark}

% --- Title and Authors ---
\title{Rank-based Optional Stopping via Structure Towers}
\author{
  Claude Code\thanks{TeX generation and structural synthesis.} \\
  \texttt{claude.ai}
  \and
  Codex (OpenAI)\thanks{Lean 4 formalization and code generation.} \\
  \texttt{openai.com}
}
\date{\today}

\begin{document}

\maketitle

\begin{abstract}
We present a structural re-interpretation of the Optional Stopping Theorem (OST) for martingales using the concept of Structure Towers and Rank Functions. By formally identifying stopping times with rank functions on a stratified probability space, we transform the classical probabilistic statement into a structural property. We formalize this correspondence in Lean 4, demonstrating that the ``stopped process'' is canonically identified with a rank-restricted evaluation. This paper details the construction of the Structure Tower induced by a stopping time, the translation of probabilistic conditions to rank properties, and the subsequent derivation of the Bounded Optional Stopping Theorem. The formalization prioritizes a ``thin wrapper'' strategy, separating the high-level rank-theoretic arguments from the underlying measure-theoretic implementation.
\end{abstract}

% --- Section 1: Introduction ---
\section{Introduction}

The Optional Stopping Theorem (OST) is a cornerstone of martingale theory, asserting that the expected value of a martingale remains invariant under random stopping times, provided certain boundedness or integrability conditions are met. Traditionally, this result is presented purely within the language of measure theory. In this paper, we propose a categorical shift in perspective: we view the stopping time not just as a measurable random variable, but as a \textit{Rank Function} that induces a \textit{Structure Tower} on the probability space.

This section provides an overview of the paper's motivation and structure. We aim to bridge the gap between abstract structuralism and concrete probability theory. By defining a canonical correspondence between stopping times and structure towers, we show that the ``stopping'' operation is equivalent to a ``minimization'' operation in a rank lattice. This abstraction simplifies the conceptual burden of the proof and aligns naturally with formal verification strategies.

The primary contributions of this work are:
\begin{itemize}
    \item \textbf{Formal Equivalence:} We establish a correspondence between Structure Towers (nested subsets satisfying a covering property) and Rank Functions (maps to $\Nat$), formalized as a ``Normal Form'' theorem.
    \item \textbf{Stopping Times as Ranks:} We prove that any stopping time $\tau$ on a filtration $\Filtration$ induces a canonical rank function $\rank$, where the event $\{\tau \le n\}$ corresponds exactly to the $n$-th layer of the associated structure tower.
    \item \textbf{Rank-based OST:} We reformulate the Bounded Optional Stopping Theorem. We prove that the expectation of a martingale is invariant under ``rank restriction'' in the structure tower.
    \item \textbf{Lean 4 Formalization:} We describe a modular formalization of these concepts, distinguishing between a ``thin wrapper'' of rank-theoretic definitions and the robust measure-theoretic backend.
\end{itemize}

The remainder of this paper is organized as follows. Section \ref{sec:discrete} establishes the discrete probability setting. Section \ref{sec:ranktower} introduces the abstract theory of Structure Towers. Section \ref{sec:stopping_ranks} maps stopping times to these towers. Section \ref{sec:ost} proves the Rank-based OST. Section \ref{sec:lean} discusses the mechanization, and Section \ref{sec:conclusion} concludes.

% --- Section 2: Discrete Setting ---
\section{Discrete Setting} \label{sec:discrete}

We restrict our attention to a discrete-time setting to focus on the structural correspondence without the distracting technicalities of continuous-time analysis.

Let $(\Omega, \mathcal{F}, \Prob)$ be a probability space. We consider a discrete time index $n \in \Nat$.
\begin{definition}[Filtration]
A filtration $\Filtration = (\mathcal{F}_n)_{n \in \Nat}$ is a sequence of sub-$\sigma$-algebras such that $\mathcal{F}_n \subseteq \mathcal{F}_{n+1} \subseteq \mathcal{F}$ for all $n$.
\end{definition}

\begin{definition}[Martingale]
A stochastic process $M = (M_n)_{n \in \Nat}$ is a martingale with respect to $\Filtration$ if:
\begin{enumerate}
    \item $M_n$ is $\mathcal{F}_n$-measurable for all $n$.
    \item $\Ex[|M_n|] < \infty$ for all $n$ (Integrability).
    \item $\CondExp{M_{n+1}}{n} = M_n$ almost surely.
\end{enumerate}
\end{definition}

\begin{definition}[Stopping Time]
A map $\tau: \Omega \to \Nat \cup \{\infty\}$ is a stopping time with respect to $\Filtration$ if for every $n \in \Nat$, the event $\{\omega \in \Omega \mid \tau(\omega) \le n\}$ belongs to $\mathcal{F}_n$.
\end{definition}

Our goal is to reinterpret these standard definitions using the geometry of ``layers'' and ``ranks.''

% --- Section 3: Minimal Structure-Tower Definitions ---
\section{Structure Towers and Rank Functions} \label{sec:ranktower}

In this section, we define the abstract machinery of Structure Towers. This theory, formalized in \texttt{RankTower.md}, provides the canonical correspondence between stratified structures and rank functions.

\begin{definition}[Structure Tower]
A \textbf{Structure Tower with Min} consists of a carrier type $X$ and a sequence of subsets (layers) $\layer{n} \subseteq X$ satisfying:
\begin{enumerate}
    \item \textbf{Covering:} $\forall x \in X, \exists n, x \in \layer{n}$.
    \item \textbf{Monotonicity:} $i \le j \implies \layer{i} \subseteq \layer{j}$.
    \item \textbf{Minimality:} There exists a function $\minLayer: X \to \Nat$ such that $\minLayer(x)$ is the smallest index $n$ where $x \in \layer{n}$.
\end{enumerate}
\end{definition}

\begin{definition}[Rank Function]
A \textbf{Rank Function} is a map $\rank: X \to \Nat \cup \{\infty\}$. We say a rank function is \textit{covering} if $\forall x, \rank(x) < \infty$, in which case it can be identified with a map $\rank: X \to \Nat$.
\end{definition}

\begin{remark}
Although the formal definition allows for infinite ranks (implemented as \texttt{WithTop Nat} in Lean), we restrict our focus in this paper to covering rank functions (corresponding to bounded or finite stopping times) to establish a clean bijection with the Structure Tower defined above.
\end{remark}

The following theorem establishes the consistency between the layer structure and the minimal layer function.

\begin{theorem}[Rank-Tower Correspondence] \label{thm:rank_tower_corr}
There is a one-to-one correspondence between Structure Towers and Covering Rank Functions ($\rank: X \to \Nat$).
\begin{enumerate}
    \item \textbf{Rank to Tower:} Given $\rank$, define $\layer{n} = \{x \mid \rank(x) \le n\}$.
    \item \textbf{Tower to Rank:} Given a tower, define $\rank(x) = \minLayer(x)$.
\end{enumerate}
These operations are inverses: the $\minLayer$ derived from the induced tower of $\rank$ is exactly $\rank$, and the layers induced by $\minLayer$ are exactly the original layers.
\end{theorem}

\begin{proof}[Sketch]
The monotonicity ensures that the set of indices containing $x$ forms a filter in $\Nat$, determined by its minimum. The formalization \texttt{RankTower.lean} proves the round-trip identity, effectively showing that the explicit data of \texttt{minLayer} in the structure definition is redundant but consistent with the layer data.
\end{proof}

% --- Section 4: Stopping Times as Ranks ---
\section{Stopping Times as Ranks} \label{sec:stopping_ranks}

We now bridge the gap between structure towers and probability theory.
Throughout this section and the next, we assume $\tau$ is a \textbf{bounded stopping time}. That is, there exists $K \in \Nat$ such that $\tau(\omega) \le K$ for all $\omega$. Consequently, $\tau$ takes values in $\Nat$, allowing us to treat it as a covering rank function $\rank_\tau: \Omega \to \Nat$.

\begin{proposition}[Stopping Time as Rank]
Let $\tau$ be a bounded stopping time. Define $\rank_\tau(\omega) = \tau(\omega)$. Then the induced structure tower $\Tower_\tau$ satisfies:
\begin{equation}
    \layer{n}(\Tower_\tau) = \{\omega \in \Omega \mid \tau(\omega) \le n\}.
\end{equation}
\end{proposition}

While $\minLayer(\omega) = \tau(\omega)$ holds by definition, the probabilistic significance lies in the following observation:

\begin{remark}[Non-triviality of Adaptedness]
The structural definition alone yields subsets $\layer{n} \subseteq \Omega$. The crucial probabilistic content is that because $\tau$ is a stopping time, we have $\layer{n}(\Tower_\tau) \in \mathcal{F}_n$. This means the geometry of the structure tower is ``adapted'' to the filtration $\Filtration$. The rank structure encodes the information flow.
\end{remark}

We define the \textbf{Rank-Stopped Process} $M^\tau$ (denoted \texttt{rankStoppedProcess} in the formalization) and the \textbf{Terminal Variable} $M_\tau$.
\begin{definition}
\begin{itemize}
    \item The process stopped at rank $n$: $M^\tau_n(\omega) := M_{\min(n, \tau(\omega))}(\omega)$.
    \item The value at the stopping time: $M_\tau(\omega) := M_{\tau(\omega)}(\omega)$.
\end{itemize}
Since $\tau$ is bounded by $K$, for any $n \ge K$, we have $M^\tau_n = M_\tau$.
\end{definition}

% --- Section 5: Rank Optional Stopping Theorem ---
\section{Rank-based Optional Stopping Theorem} \label{sec:ost}

We now present the main result: the Bounded Optional Stopping Theorem formulated via rank structures.

\begin{theorem}[Rank OST - Bounded Case] \label{thm:rank_ost}
Let $M$ be a martingale and $\tau$ a bounded stopping time (bound $K$). Then for any time $n$:
\[ \Ex[M^\tau_n] = \Ex[M_0]. \]
Consequently, taking $n=K$, we have $\Ex[M_{\tau}] = \Ex[M_0]$.
\end{theorem}

\begin{proof}[Proof Sketch]
The proof relies on two main steps:
\begin{enumerate}
    \item \textbf{Martingale Property Preservation:} We first show that the stopped process $M^\tau$ is a martingale. The key calculation is the conditional expectation of the increment. Using the fact that $\{\tau \le n\} \in \mathcal{F}_n$:
    \begin{align*}
        \CondExp{M^\tau_{n+1}}{n} &= \CondExp{M_{\tau \land (n+1)}}{n} \\
        &= \CondExp{\Ind{\tau \le n} M_\tau + \Ind{\tau > n} M_{n+1}}{n} \\
        &= \Ind{\tau \le n} M_\tau + \Ind{\tau > n} \CondExp{M_{n+1}}{n} \\
        &= \Ind{\tau \le n} M_\tau + \Ind{\tau > n} M_n \\
        &= M_{\tau \land n} = M^\tau_n.
    \end{align*}
    Integrability of $M^\tau_n$ follows from the integrability of $M$ and the boundedness of $\tau$ (it is a finite sum of integrable variables).
    
    \item \textbf{Expectation Conservation:} Since $M^\tau$ is a martingale, $\Ex[M^\tau_n] = \Ex[M^\tau_0]$. Since $\tau \ge 0$, $M^\tau_0 = M_0$. Thus $\Ex[M^\tau_n] = \Ex[M_0]$.
\end{enumerate}
\end{proof}

\begin{corollary}[Rank Expectation Invariance]
For a bounded rank function $\tau$, the expectation is invariant under the choice of layer $n$ in the rank tower.
\end{corollary}

% --- Section 6: Lean Mechanization ---
\section{Lean Mechanization} \label{sec:lean}

The formalization of these results in Lean 4 adopts a ``thin wrapper'' strategy.

\subsection{Architecture}
The project is divided into distinct layers:
\begin{enumerate}
    \item \textbf{\texttt{RankTower.lean}:} Pure order theory. Defines \texttt{StructureTowerWithMin} and proves the correspondence with rank functions.
    \item \textbf{\texttt{StoppingTime\_C.lean}:} Interface layer. Defines \texttt{stoppingTimeAsRank} and proves that the layer sets equal the stopping sets.
    \item \textbf{\texttt{Martingale\_RankStructure.lean}:} Computational core. It wraps existing Mathlib lemmas into rank terminology.
    \item \textbf{\texttt{OptionalStopping\_RankTheory.lean}:} High-level theory. States the OST purely in terms of ranks.
\end{enumerate}

\subsection{Benefits of the Thin Wrapper}
The ``thin wrapper'' approach simplifies the application of standard probability theorems. For example, to prove the conservation of expectation for the rank-stopped process, we do not need to reconstruct the integral theory. Since \texttt{rankStoppedProcess M tau} is definitionally equal to the standard \texttt{Martingale.stoppedProcess}, we can directly apply Mathlib's \texttt{integral\_condExp} and \texttt{martingale\_of\_bdd}.

\begin{verbatim}
-- In Lean, the proof becomes trivial due to definitional equality:
theorem rankOptionalStopping_bounded ... :
  Integral (rankStoppedProcess M tau n) = Integral (M 0) := by
  -- The definition unfolds to the standard stopped process
  exact Martingale.stopped_process_integral_eq ...
\end{verbatim}

This allows the user to reason at the high level of Rank Towers while the backend transparently handles the measure-theoretic details.

% --- Section 7: Discussion ---
\section{Discussion} \label{sec:conclusion}

We have presented a structural approach to the Optional Stopping Theorem. By identifying stopping times with rank functions, we provided a geometric interpretation of the OST as an invariance property across the layers of a structure tower. While we focused on the bounded case, the correspondence paves the way for treating unbounded stopping times via limits in the structure tower.

% --- Appendices ---
\appendix

\section{RankTower Normal Form (Lean 4)}
The following snippet from \texttt{RankTower.md} illustrates the definition of the structure tower. Symbols have been transliterated to ASCII.

\begin{verbatim}
structure StructureTowerWithMin where
  carrier : Type*
  layer : Nat -> Set carrier
  covering : forall x : carrier, exists i : Nat, x \in layer i
  monotone : forall {i j : Nat}, i <= j -> layer i \subseteq layer j
  minLayer : carrier -> Nat
  minLayer_mem : forall x, x \in layer (minLayer x)
  minLayer_minimal : forall x i, x \in layer i -> minLayer x <= i
\end{verbatim}

% --- Bibliography ---
\begin{thebibliography}{9}

\bibitem{williams1991}
D. Williams,
\textit{Probability with Martingales},
Cambridge University Press, 1991.

\bibitem{mathlib}
The Mathlib Community,
\textit{The Lean Mathematical Library},
arXiv:1910.09336, 2019.

\end{thebibliography}

\end{document}